\section{Experimental Evaluation}

\subsection{Setup}

\textbf{Model and Hardware.} Llama-3-8B (32 layers, $d_{\text{model}}\!=\!4096$, $n_{\text{heads}}\!=\!32$, $n_{\text{kv\_heads}}\!=\!8$, $d_{\text{head}}\!=\!128$), context length 16,400, 128 decode steps, batch size 64.
GPU: B100 (989 TFLOPS FP16, 3.35 TB/s HBM).
PIM: 2048 banks (8 stacks $\times$ 256), 2KB row buffer, PE 16B/cycle at 2GHz, 32 GB/s per-bank bandwidth.
Asymmetric quantization: $z_Q\!=\!1.0$, $z_K\!=\!2.0$.

\textbf{Simulator.} Cycle-accurate C++ simulator extending LLMSimulator~\cite{duplex} with configurable near-bank PIM units.
Reports per-stage timing (Q@K, softmax, score@V) and per-bank component times (rowbuffer vs. PE).
All PIM latencies verified against analytical cycle counts (ratio 1.00$\times$).

\textbf{Configurations.} Six configurations evaluated:
\begin{enumerate}
\item \textbf{FP16 GPU}: All GPU, float16 KV.
\item \textbf{FP16 PIM}: All PIM, float16 KV (reference).
\item \textbf{2-bit GPU}: 2-bit KV, all GPU.
\item \textbf{2-bit Hybrid} (\pimattn{}): Q@K on PIM, softmax+score@V on GPU.
\item \textbf{2-bit All-PIM}: Both Q@K and score@V on PIM.
\item \textbf{2-bit Dequant}: Dequantize to FP16 in PIM before GEMV.
\end{enumerate}

\subsection{Per-Step Attention Latency}

Table~\ref{tab:main} presents per-step attention latency (16 seqs/DP group).

\begin{table}[t]
\caption{Per-step attention latency ($\mu s$).}
\label{tab:main}
\centering
\begin{tabular}{lccccc}
\toprule
\textbf{Config} & \textbf{Q@K} & \textbf{S@V} & \textbf{Gen} & \textbf{pim\_rb} & \textbf{pim\_pe} \\
\midrule
FP16 GPU   & 97 & --- & 95  & --- & --- \\
FP16 PIM   & 66 & 66  & 132 & 131 & 66  \\
2-bit GPU  & 48 & --- & 48  & --- & --- \\
2-bit Hyb. & 33 & 34  & 67  & 33  & 17  \\
2-bit All  & 33 & 33  & 66  & 66  & 33  \\
2-bit Deq. & 66 & 66  & 132 & 66  & 99  \\
\bottomrule
\end{tabular}
\end{table}

\textbf{Key findings:}
\begin{itemize}
\item \pimattn{} (2-bit Hybrid) achieves $2.1\mu s$/seq Q@K vs. $3.0\mu s$/seq for 2-bit GPU --- $1.4\times$ speedup. Against FP16 GPU ($6.0\mu s$/seq), the advantage is $2.9\times$.
\item Dequantize-then-compute is inefficient: $3\times$ PE increase (99 vs. 33 $\mu s$), doubling Q@K latency. This validates direct 2-bit compute.
\item Row-buffer time dominates: $pim\_rb$ accounts for 50--66\% of latency, indicating memory-bandwidth-bound execution at 2048 banks.
\item Score@V benefits from GPU: at 34$\mu s$ (GPU) vs. 33$\mu s$ (PIM), GPU float16 throughput offsets data movement overhead.
\end{itemize}

\subsection{PIM Cycle Verification}

For FP16 PIM at $N_{\text{banks}}\!=\!2048$: per-GEMV data = 4.26MB, per-bank = 2082B, requiring 2 row-buffer fills.
Row-buffer time: $2 \times 64\text{ns} = 128\text{ns}$ per GEMV.
PE time: $2082\text{B}/16 \times 0.5\text{ns} = 65.1\text{ns}$ per GEMV.
Per-step (Q@K + score@V, $\times 4$ group, $\times 8$ heads, $\times 16$ seqs): RB $131.1\mu s$, PE $66.6\mu s$.
\textbf{Reported: $131.1\mu s$ and $66.5\mu s$ (ratio 1.00$\times$), confirming simulator fidelity.}

\subsection{Asymmetric Quantization Overhead}

With asymmetric quantization, reduction terms add PE time proportional to $(M+N)\!\cdot\!K$.
For decode mode ($M\!=\!1, N\!=\!16528$), this matches the main GEMV data volume, resulting in $\sim\!1.8\times$ PE time increase.
In prefill mode with larger $M$, the relative overhead diminishes.

\subsection{End-to-End Throughput}

Table~\ref{tab:e2e} reports total latency for 128 steps $\times$ 32 layers.

\begin{table}[t]
\caption{End-to-end latency ($\mu s$) for 128 steps $\times$ 32 layers.}
\label{tab:e2e}
\centering
\begin{tabular}{lcc}
\toprule
\textbf{Config} & \textbf{Total} & \textbf{Per Layer} \\
\midrule
FP16 GPU   & 5,623  & 176 \\
FP16 PIM   & 6,788  & 212 \\
2-bit GPU  & 2,769  & 87  \\
2-bit Hyb. & 3,460  & 108 \\
2-bit All  & 3,443  & 108 \\
2-bit Deq. & 5,555  & 174 \\
\bottomrule
\end{tabular}
\end{table}

2-bit GPU achieves the lowest latency (2,769$\mu s$) due to reduced memory traffic combined with GPU throughput.
\pimattn{} (3,460$\mu s$) is $1.25\times$ slower than 2-bit GPU but offers the advantage of keeping KV cache in PIM memory, enabling larger batch sizes and longer contexts exceeding GPU memory capacity.
